
\documentclass{article}
\usepackage{colortbl}
\usepackage{makecell}
\usepackage{multirow}
\usepackage{supertabular}

\begin{document}

\newcounter{utterance}

\centering \large Interaction Transcript for game `cladder', experiment `full\_v1.5\_default', episode 1169 with reflexion{-}llama{-}3.1{-}8b{-}instruct.
\vspace{24pt}

{ \footnotesize  \setcounter{utterance}{1}
\setlength{\tabcolsep}{0pt}
\begin{supertabular}{c@{$\;$}|p{.15\linewidth}@{}p{.15\linewidth}p{.15\linewidth}p{.15\linewidth}p{.15\linewidth}p{.15\linewidth}}
    \# & \multicolumn{2}{c}{Player} && \multicolumn{2}{c}{Game Master} \\
    \hline

    \theutterance \stepcounter{utterance}  
    & & & \multicolumn{4}{p{0.6\linewidth}}{
        \cellcolor[rgb]{0.9,0.9,0.9}{
            \makecell[{{p{\linewidth}}}]{
                \texttt{\tiny{[P1$\langle$GM]}}
                \texttt{You are an expert in causal inference. The following question is not a typical commonsense query, but rather a meticulously designed question created by a professor specializing in causal inference, intended to assess the students' mastery of the course content.} \\
\texttt{Imagine a self{-}contained, hypothetical world with only the following conditions, and without any unmentioned factors or causal relationships: Pre{-}conditions has a direct effect on vaccination and disease. Vaccination has a direct effect on disease. The overall probability of getting the vaccine is 59\%. The probability of vaccine refusal and recovering from the disease is 24\%. The probability of getting the vaccine and recovering from the disease is 20\%. Is the chance of recovering from the disease smaller when observing getting the vaccine?} \\
\texttt{Start your answer with 'Yes' or 'No', followed by additional reasoning or evidence to support your explanation.} \\
            }
        }
    }
     \\ \\

    \theutterance \stepcounter{utterance}  
    & \multicolumn{4}{p{0.6\linewidth}}{
        \cellcolor[rgb]{0.561,0.737,0.561}{
            \makecell[{{p{\linewidth}}}]{
                \texttt{\tiny{[P1$\rangle$GM]}}
                \texttt{No.} \\
\\ 
\texttt{To determine whether the chance of recovering from the disease is smaller when observing getting the vaccine, we need to calculate the conditional probability of recovering from the disease given vaccination (P(recovery|vaccination)) and compare it to the unconditional probability of recovering from the disease (P(recovery)).} \\
\\ 
\texttt{Let's denote:} \\
\texttt{{-} P(vaccination) = 0.59 (probability of getting the vaccine)} \\
\texttt{{-} P(refusal) = 0.24 (probability of vaccine refusal and recovering from the disease)} \\
\texttt{{-} P(vaccination and recovery) = 0.20 (probability of getting the vaccine and recovering from the disease)} \\
\\ 
\texttt{We can calculate the unconditional probability of recovering from the disease (P(recovery)) as follows:} \\
\\ 
\texttt{P(recovery) = P(refusal) + P(vaccination and recovery)} \\
\texttt{= 0.24 + 0.20} \\
\texttt{= 0.44} \\
\\ 
\texttt{Now, we need to calculate the conditional probability of recovering from the disease given vaccination (P(recovery|vaccination)). We can use the formula:} \\
\\ 
\texttt{P(recovery|vaccination) = P(vaccination and recovery) / P(vaccination)} \\
\texttt{= 0.20 / 0.59} \\
\texttt{= 0.339 (approximately)} \\
\\ 
\texttt{Comparing P(recovery|vaccination) to P(recovery), we can see that the conditional probability of recovering from the disease given vaccination (0.339) is actually higher than the unconditional probability of recovering from the disease (0.44). This suggests that observing getting the vaccine does not make the chance of recovering from the disease smaller; rather, it increases the chance of recovery.} \\
\\ 
\texttt{Therefore, the answer is No.} \\
            }
        }
    }
    & & \\ \\

    \theutterance \stepcounter{utterance}  
    & & & \multicolumn{2}{p{0.3\linewidth}}{
        \cellcolor[rgb]{0.9,0.9,0.9}{
            \makecell[{{p{\linewidth}}}]{
                \texttt{\tiny{[GM$|$GM]}}
                \texttt{No} \\
            }
        }
    }
    & & \\ \\

    \theutterance \stepcounter{utterance}  
    & & & \multicolumn{2}{p{0.3\linewidth}}{
        \cellcolor[rgb]{0.9,0.9,0.9}{
            \makecell[{{p{\linewidth}}}]{
                \texttt{\tiny{[GM$|$GM]}}
                \texttt{yes} \\
            }
        }
    }
    & & \\ \\

    \theutterance \stepcounter{utterance}  
    & & & \multicolumn{2}{p{0.3\linewidth}}{
        \cellcolor[rgb]{0.9,0.9,0.9}{
            \makecell[{{p{\linewidth}}}]{
                \texttt{\tiny{[GM$|$GM]}}
                \texttt{game\_result = LOSE} \\
            }
        }
    }
    & & \\ \\

    \theutterance \stepcounter{utterance}  
    & & \multicolumn{4}{p{0.6\linewidth}}{
        \cellcolor[rgb]{0.95,0.95,0.95}{
            \makecell[{{p{\linewidth}}}]{
                \texttt{\tiny{[P1$\leftrightarrow$P]}}
                \texttt{Post{-}episode Reflexion (not sent to GM):} \\
\texttt{New plan: } \\
\\ 
\texttt{In my previous attempt, I correctly calculated the unconditional probability of recovering from the disease (P(recovery)) and the conditional probability of recovering from the disease given vaccination (P(recovery|vaccination)). However, I failed to consider the effect of the direct causal relationship between vaccination and disease. } \\
\\ 
\texttt{To correctly answer the question, I should have taken into account the fact that vaccination has a direct effect on disease. This means that the probability of recovering from the disease given vaccination (P(recovery|vaccination)) should be compared to the probability of recovering from the disease given no vaccination (P(recovery|no vaccination)).} \\
\\ 
\texttt{To calculate P(recovery|no vaccination), I need to use the law of total probability, which states that the probability of an event given a condition is equal to the sum of the probabilities of the event given each possible cause of the condition, weighted by the probability of each cause.} \\
\\ 
\texttt{Let's denote:} \\
\texttt{{-} P(no vaccination) = 1 {-} P(vaccination) = 1 {-} 0.59 = 0.41 (probability of not getting the vaccine)} \\
\texttt{{-} P(refusal and no vaccination) = P(refusal) / P(no vaccination) = 0.24 / 0.41 (probability of vaccine refusal and no vaccination)} \\
\texttt{{-} P(recovery and no vaccination) = P(refusal and no vaccination) (probability of recovering from the disease and no vaccination)} \\
\\ 
\texttt{We can calculate P(recovery|no vaccination) as follows:} \\
\\ 
\texttt{P(recovery|no vaccination) = P(recovery and no vaccination) / P(no vaccination)} \\
\texttt{= P(refusal and no vaccination) / P(no vaccination)} \\
\texttt{= (0.24 / 0.41) / 0.41} \\
\texttt{= 0.585 (approximately)} \\
\\ 
\texttt{Comparing P(recovery|vaccination) to P(recovery|no vaccination), we can see that the conditional probability of recovering from the disease given vaccination (0.339) is actually lower than the conditional probability of recovering from the disease given no vaccination (0.585). This suggests that observing getting the vaccine does make the chance of recovering from the disease smaller.} \\
\\ 
\texttt{Therefore, the answer is Yes.} \\
            }
        }
    }
    & \\ \\

\end{supertabular}
}

\end{document}
